\item  \textbf{Diets.}  Developers of diets like to advertise that their diets produce significant weight loss results.  A new diet has been developed and several studies were done to investigate the effectiveness of the diet. In each study the average number of pounds lost was of interest, and researchers would like to show that dieters lost a significant amount of weight.
\begin{enumerate}
\item Identify $\mu$ in this scenario. \textbf{(Choose One)}
\begin{itemize}
	\item Average number of pounds lost by all people in the sample.
	\item True average weight of all people in the population.
	\item True average number of pounds lost by all people in the population.
	\item Average weight of all people in the sample.
\end{itemize}
\newpage
\item  Identify the appropriate null and the alternative hypotheses to test the researchers' claim.
\begin{center}
\vspace{.5cm}
$H_0:\underline{\hspace{3cm}} \hspace{4cm} H_a:\underline{\hspace{3cm}}$
\end{center}
\item[]
\item Suppose that an average weight loss of 1 pound (${\overline x} = 1$) and standard deviation of $s = 7$ pounds were found in each of the following scenarios. Further, assume weight loss is normally distributed. For each case, compute the test statistic, find the approximate $p$-value and make a decision about the null hypothesis using an $\alpha = 0.05$ level of significance. Round $t$-statistics to two decimals.

\begin{enumerate}
\item \textbf{\underline{Scenario 1:}} a sample of size $n = 10$\\
\begin{enumerate}
\vspace{0.25cm}
\item $t = \underline{\hspace{4cm}}$ and df = \underline{\hspace{4cm}}
\item The $p$-value corresponding to the correct test statistic is 0.331. This $p$-value suggests $\underline{\hspace{4cm}}$ against the null hypothesis. \textbf{(Choose One)}
\begin{itemize}
\item Little or no evidence
\item Weak evidence
\item Moderately strong evidence
\item Strong evidence
\item Extremely strong evidence
\end{itemize}

\item \underline{Decision:}
\begin{itemize}
\item Reject the null hypothesis because the $p$-value $\leq$ $\alpha$.
\item Reject the null hypothesis because the $p$-value $>$ $\alpha$.
\item Fail to reject the null hypothesis because the $p$-value $\leq$ $\alpha$.
\item Fail to reject the null hypothesis because the $p$-value $>$ $\alpha$.
\end{itemize}
\end{enumerate}
\item[]
%\newpage
\item \textbf{\underline{Scenario 2:}} a sample of size $n = 120$\\
\begin{enumerate}
\vspace{0.25cm}
\item $t = \underline{\hspace{4cm}}$  and df = \underline{\hspace{4cm}}
\item The $p$-value corresponding to the correct test statistic is 0.0601. This $p$-value suggests $\underline{\hspace{4cm}}$ against the null hypothesis. \textbf{(Choose One)}
\begin{itemize}
\item Little or no evidence
\item Weak evidence
\item Moderately strong evidence
\item Strong evidence
\item Extremely strong evidence
\end{itemize}

\item \underline{Decision:}
\begin{itemize}
\item Reject the null hypothesis because the $p$-value $\leq$ $\alpha$.
\item Reject the null hypothesis because the $p$-value $>$ $\alpha$.
\item Fail to reject the null hypothesis because the $p$-value $\leq$ $\alpha$.
\item Fail to reject the null hypothesis because the $p$-value $>$ $\alpha$.
\end{itemize}
\end{enumerate}
\newpage
%\item[]
\item \textbf{\underline{Scenario 3:}} a sample of size $n = 1050$\\
\begin{enumerate}
\vspace{0.25cm}
\item $t = \underline{\hspace{4cm}}$  and df = \underline{\hspace{4cm}}
\item  The $p$-value corresponding to the correct test statistic is 0.000002. This $p$-value suggests $\underline{\hspace{4cm}}$ against the null hypothesis. \textbf{(Choose One)}
\begin{itemize}
\item Little or no evidence
\item Weak evidence
\item Moderately strong evidence
\item Strong evidence
\item Extremely strong evidence
\end{itemize}

\item \underline{Decision:}
\begin{itemize}
\item Reject the null hypothesis because the $p$-value $\leq$ $\alpha$.
\item Reject the null hypothesis because the $p$-value $>$ $\alpha$.
\item Fail to reject the null hypothesis because the $p$-value $\leq$ $\alpha$.
\item Fail to reject the null hypothesis because the $p$-value $>$ $\alpha$.
\end{itemize}
%\end{enumerate}
%\end{enumerate}
\item[]
\item  (\textbf{FREE RESPONSE}) A weight loss of one pound is found to be statistically significant. The researchers share their results with the public claiming that their diet produces meaningful, significant weight loss. Do you agree? Explain your answer.
\end{enumerate}
%\end{enumerate}

%\vspace{1.5cm}

%\item Based on the results of your three tests, what can you say about the effect of sample size on producing statistically significant results? Select the all the true responses.
%\begin{enumerate}
%\item Increasing the sample size can make even a small effect practically significant.
%\item Small sample sizes increase the probability of making a type I error.
%\item Increasing the sample size will increase the size of an effect.
%\item Increasing the sample size can make even a small effect statistically significant.
%\item At large sample sizes, all statistically significant effects are practically significant.
%\end{enumerate}•
%Solutions to the above two questions will be provided.
\end{enumerate}
\end{enumerate}